% RETYPE direct from scan, replaces mathematically wrong earlier source draft
% Source: Gauss, Werke, Band II, "Recension der Untersuchungen ueber die
%         Eigenschaften der positiven ternaeren quadratischen Formen
%         von Ludwig August Seeber" (1831), p. 192--193 of the original.
% Scan: werkecarlf02gausrich.pdf p. 203 (= original Werke II p. 193).
% Audit reference: GAUSS_BANDS_I_II_III_FIDELITY.md (Band II Seeber identity,
% multiple prime-dropping errors + sign/coefficient corruption).
\documentclass[11pt]{article}
\usepackage[a4paper,margin=2.5cm]{geometry}
\usepackage{amsmath,amssymb}
\usepackage[T1]{fontenc}
\usepackage{lmodern}
\usepackage{microtype}

\title{Recension der Seeber'schen Untersuchungen\\[2pt]
\large (Gauss, Werke, Band II, p.\,193: die identischen Gleichungen)}
\author{C.\,F.\ Gauss\\[2pt]\normalsize Re-typeset from scan}
\date{}

\setlength{\parskip}{4pt}
\setlength{\parindent}{1.4em}

\begin{document}
\maketitle

\section*{Vorbemerkung}

Die positive ternäre quadratische Form
\[
  axx+byy+czz+2dyz+2bzx+2cxy
\]
(Determinante $=-D$) sei reducirt im Sinne von Seeber; mit $a',b',c'$ werden
diejenigen Coefficienten bezeichnet, die in den Vorgleichungen ($\dots$) als
$2a', 2b', 2c'$ vorkommen.

\section*{I.\ Fall: keine der Zahlen $a',b',c'$ negativ}

Setze man
\[
  b-2a'=d,\quad c-2b'=e,\quad a-2c'=f,
\]
\[
  c+2a'=g,\quad a+2b'=h,\quad b+2c'=i,
\]
\emph{wo aus der Definition der reducirten positiven Formen sogleich folgt,}
dass wenn $axx+byy+czz+2dyz+2bzx+2cxy$ eine solche ist, keine jener sechs
Zahlen negativ ist. Bezeichnet man nun den (negativen) Determinanten der Form
durch $-D$, so hat man, wie man sich durch die Entwickelung leicht überzeugt,
die identische Gleichung
\[
\boxed{\;2D-abc \;=\; a\,a'\,d+b\,b'\,e+c\,c'\,f+a'\,h\,i+b'\,g\,i+c'\,g\,h+g\,h\,i\;}
\]
in welcher keines der sieben Glieder zur Rechten negativ sein kann, und folglich
$abc$ nicht grösser als $2D$. Dasselbe folgt auf gleiche Weise aus der
identischen Gleichung
\[
\boxed{\;2D-abc \;=\; a\,a'\,g+b\,b'\,h+c\,c'\,i+a'\,e\,f+b'\,d\,f+c'\,d\,e+d\,e\,f\;}
\]

\section*{II.\ Fall: keine der Zahlen $a',b',c'$ positiv}

Setze man
\[
  b+2a'=d,\quad c+2b'=e,\quad a+2c'=f,
\]
\[
  c+2a'=g,\quad a+2b'=h,\quad b+2c'=i,
\]
\[
  b+c+2a'+2b'+2c'=k,
\]
\[
  a+c+2a'+2b'+2c'=l,
\]
\[
  a+b+2a'+2b'+2c'=m,
\]
und den Determinanten der Form wie vorhin $=-D$. Vermöge der Definition der
reducirten positiven Formen wird keine der neun Zahlen $d,e,f,g,h,i,k,l,m$
negativ sein können, und so ergibt sich aus der identischen Gleichung
\begin{multline*}
6D-3abc \;=\; -a\,a'(d+2k)-b\,b'(e+2l)-c\,c'(f+2m)\\
{}-a'\,h\,i-b'\,g\,i-c'\,g\,h+d\,e\,f+2\,g\,h\,i
\end{multline*}
in welcher, weil $a',b',c'$ nicht positiv, sondern negativ oder Null sind, alle
Glieder zur Rechten positiv oder Null werden, dass also $3abc$ nicht grösser als
$6D$, oder $abc$ nicht grösser als $2D$ sein kann. Dasselbe folgt eben so aus
der identischen Gleichung
\begin{multline*}
6D-3abc \;=\; -a\,a'(g+2k)-b\,b'(h+2l)-c\,c'(i+2m)\\
{}-a'\,e\,f-b'\,d\,f-c'\,d\,e+2\,d\,e\,f+g\,h\,i
\end{multline*}

\section*{Vereinfachte unsymmetrische Gleichung}

\emph{Beide Gleichungen sind symmetrisch. Verzichtet man auf völlige Symmetrie,
so ist der Beweis mit einer noch geringeren Anzahl von Gliedern zu führen, z.\,B.\
durch die identische Gleichung}
\[
\boxed{\;8D-4abc \;=\; -2\,a\,a'(g+k)-2\,b\,b'(e+l)-4\,c\,c'\,m+(c+e)\,d\,f+(c+g)\,h\,i\;}
\]

\bigskip
\hrule\medskip
\noindent\footnotesize Audit notes: The reader (earlier source draft) version of this passage
had at least the following errors, all corrected here against the scan:
(1) $2D-abc$ was misread as $\tfrac12 D=abc$ in Eq.\,A and $4D=abc$ in Eq.\,B
    (sign and operator flipped, coefficient corrupted);
(2) the primes on $a,b,c$ in $a\,a'd$, $b\,b'e$, $c\,c'f$ etc.\ were systematically
    dropped to $aad$, $bbe$, $ccf$;
(3) one term was rewritten as the wrong letters ($a'hi\to dhf$ etc.);
(4) the Case II definitions had primes dropped and were mis-grouped, producing
    the nonsense $m=a+b+c+2a+2b+2c=3a+3b+3c$;
(5) Equation~(25) $8D-4abc=\dots$ was reduced to a mangled $4abc=\dots$ with
    further dropped primes. All restored here.

\end{document}
